\documentclass[../main.tex]{subfiles}

\begin{document}

\chapter{Implementation}\label{ch:implementation}

This chapter describes the implementation of the proposed monitoring tool and 
discusses the main challenges encountered during development. 
It follows the same structure as Chapter~\ref{ch:system-design} by examining 
each system layer and describing its technical implementation. 

\section{Data Aggregation Layer}\label{sec:impl-aggregation}

As described in Section~\ref{sec:data-aggregation}, the data aggregation layer 
uses log shippers to collect and parse log entries and static configuration 
files, as well as a custom interceptor that operates within the Shibboleth IdP. 
The log shipping stack consists of Filebeat for reading the source files and 
Logstash for preprocessing and transmitting the collected data. 
Both components are deployed on the IdP side. 

A key aspect of the implementation is the handling of audit log entries. 
Instead of parsing the audit log according to a fixed format, Logstash processes 
the audit fields as key-value pairs. 
This provides the technical basis for adapting to changes in the configured audit 
format and enables the monitoring system to handle the dynamic structure of the 
audit logs. 

The monitoring tool handles these dynamic changes by maintaining a field registry. 
The registry periodically retrieves the most recent audit entries stored in the 
database and extracts the names of all fields that are currently present. 
Using only recent entries ensures that fields no longer written by the IdP are 
eventually removed from the registry. 
Both the detection rules and the correlation strategies query the registry before 
they are applied, which is how the withdrawal described in 
Section~\ref{subsec:detection} is realised in practice. 

A particular difficulty of this mechanism is that a missing field does not manifest 
itself as an error. 
If the audit format of a deployment does not include a field that a rule depends on, 
no exception is raised and no entry is written to any log. 
The information is simply absent, and a rule evaluating it would silently operate on 
incomplete data. 
The registry therefore does not only serve adaptability, but also makes this absence 
explicit and detectable. 

The static configuration files are collected by the same log shipping stack, but 
processed differently. 
Unlike the audit log, which is written continuously and line by line, 
\textit{relying-party.xml} and the metadata documents change rarely and are read 
as a whole. 
Filebeat is therefore configured to transmit the complete file whenever a change 
is detected, and Logstash forwards it without field level parsing. 
The interpretation of the XML structure takes place in the monitoring tool itself, 
where the relevant security settings are extracted and stored per service provider. 

Unlike the log shippers, the interceptor is a custom component developed as part of 
this work. 
It is registered as a Spring bean within the Shibboleth IdP and activated through the 
\texttt{postAuthenticationFlows} property of the corresponding profile configuration. 
The interceptor therefore integrates into the IdP without any modification of its 
source code, which is how the third requirement of Section~\ref{sec:requirements} is met. 

The interceptor extends the \texttt{AbstractProfileAction} API of the Shibboleth IdP and 
operates on the \texttt{ProfileRequestContext}, which represents the current state of the 
profile execution~\cite{AbstractProfileAction, ProfileRequestContext}. 
From this context it accesses the child contexts that carry the information relevant to 
the monitoring purpose. 
The \texttt{RelyingPartyContext} identifies the SP for which the authentication is 
performed, the \texttt{SubjectContext} provides the authenticated subject, and the 
\texttt{InboundMessageContext} gives access to the raw SAML request as it was 
received~\cite{RelyingPartyContext, SubjectContext, InOutOperationContext}. 
Further contexts, among them the \texttt{AuthenticationContext} and the 
\texttt{SessionContext}, supply information about the authentication flow that was used 
and the session state at the time of the request~\cite{AuthenticationContext, SessionContext}. 

The information obtained from these contexts is extracted, parsed, and combined into a 
single payload. 
Sensitive identifiers are pseudonymised at this point, before the payload leaves the IdP, 
as described in Section~\ref{sec:trust}. 
The transmission of the payload is covered in Section~\ref{sec:impl-transport}. 

\section{Transport Layer}\label{sec:impl-transport}

- HMAC-Calculation und -Examination, 
- Endpoint with the filter, that validates the signature, 
- the pseudonymization with the shared Salt
- asynchronous transmission

\section{Correlation and Detection Layer}\label{sec:impl-processing}

- ElasticSearch stores all the collected data from the data collection layer

\section{Visualization Layer}\label{sec:impl-visualization}

\end{document}
